\documentclass{article}
\usepackage[utf8]{inputenc}
\usepackage{amsmath,amssymb,amsthm}
\usepackage{graphicx}
\usepackage{booktabs}
\usepackage{hyperref}
\usepackage{algorithm}
\usepackage{algorithmic}
\usepackage{subcaption}
\usepackage[margin=1in]{geometry}

\newtheorem{definition}{Definition}
\newtheorem{proposition}{Proposition}
\newtheorem{theorem}{Theorem}

\title{What Should a GNN Ask? Task-Aware Active Edge Querying\\on Partially Observed Graphs}
\author{Pragunie Aditya}
\date{}

\begin{document}
\maketitle

\begin{abstract}
Graph neural networks propagate information along edges, yet real-world graphs are often incomplete.
When a target node is disconnected from all labelled nodes, no classifier can help it---the information is simply unreachable.
We introduce \emph{active edge querying}: given a budget of~$B$ oracle calls that reveal whether a candidate edge exists, choose queries that maximise node classification accuracy on a designated target set.
While existing active learning on graphs asks ``which node should I label?'', we ask ``which edge should I query?''---a complementary and previously unstudied question.
Our \emph{task-aware} acquisition strategy scores each candidate edge by the product of its link probability (estimated by a VGAE) and its expected improvement on target classification (estimated via label propagation).
On Cora and CiteSeer at 30\% edge observation, the task-aware strategy achieves $3\text{--}5\times$ higher accuracy gain per query than link-prediction heuristics and classical active-learning baselines (uncertainty sampling, QBC, information gain), despite finding fewer real edges.
We introduce a failure diagnostic that classifies target errors as structural, informational, or intrinsic, and show that task-aware querying eliminates structural failures first.
Our results hold across GCN, GAT, and GraphSAGE evaluators, and we provide evidence that the value function exhibits approximate submodularity, yielding a greedy approximation guarantee.
\end{abstract}


\section{Introduction}
\label{sec:intro}

Graph neural networks (GNNs) learn node representations by aggregating information along edges~\cite{kipf2017semi,hamilton2017inductive,velickovic2018graph}.
Their success depends critically on graph completeness: if a node's neighbourhood is missing, the GNN receives no useful signal for that node.
In practice, graphs are often partially observed---social networks miss unobserved friendships, citation networks lack references, and biological interaction networks contain untested pairs.

The standard response to incomplete graphs falls into two camps.
\emph{Link prediction} methods~\cite{zhang2018link,kipf2016variational} estimate the full adjacency from observed structure, but they operate offline, predict all edges at once, and ignore the downstream task.
\emph{Graph structure learning} (GSL) methods such as Pro-GNN~\cite{jin2020graph} and IDGL~\cite{chen2020iterative} jointly learn a graph and a GNN, but they hallucinate edges from features without ground-truth verification and may introduce false positives that degrade downstream performance.

Active learning on graphs~\cite{cai2017active,gao2018active,wu2019active} addresses the annotation bottleneck, but exclusively through \emph{node label} queries.
A complementary question is equally natural and has been overlooked:
\begin{quote}
\emph{Given a budget of $B$ oracle calls, which edges should be queried to maximise classification accuracy on a set of target nodes?}
\end{quote}

We formalise this as the \emph{active edge querying} problem (Section~\ref{sec:problem}).
The key insight is that not all edges are equally useful for a downstream task.
An edge connecting two nodes in an already-dense, well-classified cluster adds nothing; an edge bridging an isolated target node to a labelled neighbourhood can flip dozens of classifications.

\paragraph{Contributions.}
\begin{enumerate}
    \item We introduce the \textbf{active edge querying} problem for node classification on partially observed graphs, distinct from active node labelling, link prediction, and graph structure learning.
    \item We propose a \textbf{task-aware acquisition strategy} that scores candidate edges by their expected improvement on target classification, combining a VGAE link model with a label-propagation value scorer.
    \item We introduce a \textbf{failure diagnostic} that classifies misclassified targets as structural, informational, or intrinsic failures, enabling budget-aware query prioritisation.
    \item Experiments on Cora and CiteSeer show the task-aware strategy achieves $3\text{--}5\times$ the accuracy gain of baselines per query, and we provide evidence of approximate submodularity.
\end{enumerate}


\section{Related Work}
\label{sec:related}

\paragraph{Active learning on graphs.}
Existing methods select nodes to label, using criteria such as graph centrality~\cite{cai2017active}, information gain~\cite{gao2018active}, or representativeness~\cite{wu2019active}.
These assume the graph is fixed and fully observed.
Our work is complementary: we query edges, not labels.

\paragraph{Link prediction.}
Heuristic methods (common neighbours, Adamic-Adar~\cite{adamic2003friends}) and learned methods (VGAE~\cite{kipf2016variational}, SEAL~\cite{zhang2018link}) estimate edge existence.
They operate without a budget, predict all edges simultaneously, and do not consider a downstream classification task.
We use link prediction as a \emph{component} (candidate filtering), not as the objective.

\paragraph{Graph structure learning.}
Pro-GNN~\cite{jin2020graph} and IDGL~\cite{chen2020iterative} jointly learn graph structure and node representations.
They improve robustness to adversarial edges but add edges without ground-truth verification.
We compare directly against GSL: adding $B$ GSL-proposed edges versus $B$ actively queried edges (Section~\ref{sec:exp_gsl}).

\paragraph{Graph rewiring.}
Methods such as SDRF~\cite{topping2022understanding} and FoSR~\cite{karhadkar2023fosr} add virtual edges to improve message passing.
These edges do not correspond to real relationships; our queries discover ground-truth edges.


\section{Problem Formulation}
\label{sec:problem}

\begin{definition}[Active Edge Querying]
Given a partially observed graph $G_\mathrm{obs} = (V, E_\mathrm{obs})$ where $E_\mathrm{obs} \subset E_\mathrm{true}$, node features $X \in \mathbb{R}^{n \times d}$, a labelled set $L \subset V$, a target set $T \subset V \setminus L$, and a budget $B$:
\begin{itemize}
    \item An \emph{oracle} $\mathcal{O}(i,j) \in \{0,1\}$ reveals whether edge $(i,j) \in E_\mathrm{true}$.
    \item Choose a sequence of queries $q_1, \ldots, q_B$ that maximises $f(T, G_\mathrm{final})$, the classification accuracy on $T$ under the final graph $G_\mathrm{final} = G_\mathrm{obs} \cup \{\text{discovered edges}\}$.
\end{itemize}
\end{definition}

\paragraph{Relationship to active node labelling.}
Active node labelling spends budget on labels; active edge querying spends budget on structure.
At low observation rates, structure matters more than additional labels (Section~\ref{sec:exp_obs}).

\paragraph{Candidate pool.}
Querying all $\binom{n}{2}$ pairs is infeasible.
We construct a candidate pool from: (i)~feature $k$NN neighbours, (ii)~2-hop neighbours in $G_\mathrm{obs}$, and (iii)~target-focused $k$NN.
Pool recall (fraction of true hidden edges in the pool) is typically 30--40\%.


\section{Method}
\label{sec:method}

Our acquisition loop has three modules operating on the observed graph $G_\mathrm{obs}$, which grows as queries succeed.

\subsection{Module A: VGAE Candidate Generator}

A Variational Graph Autoencoder~\cite{kipf2016variational} with a 2-layer GCN encoder learns node embeddings $z_i$.
Link probability is $P(e_{ij}) = \sigma(z_i^\top z_j)$.
We use this to (i)~rank the candidate pool and (ii)~estimate the probability that each query will succeed (find a real edge).

\subsection{Module B: Task-Aware Value Scorer}

For each candidate edge $(i,j)$, we estimate the expected accuracy improvement on targets:
\begin{equation}
    \mathrm{score}(i,j) = P(e_{ij}) \cdot \mathbb{E}[\Delta \mathrm{acc}(T) \mid e_{ij} \text{ added}]
\end{equation}
The expected accuracy change is estimated via label propagation re-solves on the affected components.
We shortlist the top candidates by link probability and evaluate only those, keeping computation tractable.

Two variants:
\begin{itemize}
    \item \textbf{LP value scorer}: re-solve damped label propagation $F = (I - \alpha S)^{-1} \alpha S_L Y_L$ with the candidate edge added.
    \item \textbf{GCN value scorer}: train a GCN once; for each candidate, forward-pass with the modified adjacency and measure confidence change on targets.
\end{itemize}

\subsection{Module C: Failure Diagnostic}
\label{sec:diagnostic}

For each target $t \in T$, we classify the failure mode:
\begin{itemize}
    \item \textbf{Structural}: $t$ is in a connected component with no labelled node. No classifier can help; edge queries are the only remedy.
    \item \textbf{Informational}: $t$ is reachable from labelled nodes but the LP belief has high entropy (normalised entropy $> 0.5$). More edges may resolve the ambiguity.
    \item \textbf{Intrinsic}: $t$ is reachable and the classifier is confident, but the prediction is wrong (or right). Additional edges are unlikely to help.
\end{itemize}
The diagnostic-guided scorer (\textsc{value\_diag}) focuses budget on structural and informational failures, skipping intrinsic ones.


\subsection{The Acquisition Loop}

\begin{algorithm}
\caption{Task-Aware Active Edge Querying}
\begin{algorithmic}[1]
\REQUIRE $G_\mathrm{obs}$, features $X$, labels $L$, targets $T$, budget $B$
\STATE Train VGAE on $G_\mathrm{obs}$ to get link probabilities
\STATE Construct candidate pool $\mathcal{P}$ from feature $k$NN $\cup$ 2-hop $\cup$ target $k$NN
\FOR{$t = 1$ to $B$}
    \STATE Shortlist top candidates by $P(e_{ij})$
    \STATE Score shortlist by $P(e_{ij}) \cdot \mathbb{E}[\Delta\mathrm{acc}(T)]$
    \STATE Query oracle: $\mathcal{O}(i^*, j^*)$
    \IF{oracle returns 1}
        \STATE $E_\mathrm{obs} \leftarrow E_\mathrm{obs} \cup \{(i^*, j^*)\}$
    \ENDIF
    \STATE Remove $(i^*, j^*)$ from $\mathcal{P}$
\ENDFOR
\RETURN $G_\mathrm{final}$
\end{algorithmic}
\end{algorithm}


\section{Experiments}
\label{sec:experiments}

\subsection{Setup}

\paragraph{Datasets.}
Cora (2,708 nodes, 5,429 edges, 7 classes) and CiteSeer (3,327 nodes, 4,732 edges, 6 classes).
We observe 30\% of edges and hide the rest.

\paragraph{Budget and targets.}
Budget $B = 500$, 200 target nodes, 20 labels per class.
Checkpoints at $\{0, 25, 50, 100, 150, 200, 300, 400, 500\}$ queries.

\paragraph{Baselines.}
\begin{itemize}
    \item \emph{Heuristic}: Random, Common Neighbours, Adamic-Adar, Max Degree
    \item \emph{Link-based}: Link probability (logistic), VGAE top-$k$
    \item \emph{Active learning}: Uncertainty sampling, Query by Committee (5 GCNs), Information Gain
    \item \emph{Task-aware}: Value (LP), Value (GCN), Value+Diag
    \item \emph{Upper bound}: Oracle greedy (sees true edges)
    \item \emph{GSL control}: Feature-$k$NN, VGAE-GSL (no oracle)
\end{itemize}

\paragraph{Metrics.}
Primary: area under the LP accuracy--vs--queries curve (LP-AUC).
Secondary: GCN accuracy at endpoints.
Statistics: paired $t$-test + Wilcoxon, Holm correction, 95\% bootstrap CI.

\paragraph{Pre-registration.}
Seeds 200--229 are reserved for confirmatory runs and never used for tuning.
The primary metric was changed from GCN to LP \emph{before} any confirmatory run, as documented in our pre-registration file.

\subsection{Main Result}
\label{sec:exp_main}

% TABLE: Main result — all strategies, LP-AUC, LP gain, reach, edges found
% Generated by analysis_paper.py
\emph{[Table to be filled from experiment 4b results]}

\subsection{Observation Rate Sweep}
\label{sec:exp_obs}

% How does the advantage change with observation rate?
\emph{[Results from experiment 4c]}

\subsection{Ablation Study}
\label{sec:exp_ablation}

We ablate the contribution of each component:
\begin{itemize}
    \item \textbf{Value vs Value-All}: Is target-awareness necessary, or does scoring all nodes work equally well?
    \item \textbf{VGAE vs Logistic}: Does the link model matter?
    \item \textbf{LP vs GCN scorer}: Do LP and GCN scorers select different edges?
    \item \textbf{Value vs Value+Diag}: Does the failure diagnostic help?
\end{itemize}

% Results from experiment 4f
\emph{[Table to be filled from experiment 4f results]}

\subsection{Active Querying vs Graph Structure Learning}
\label{sec:exp_gsl}

% Feature-kNN and VGAE-GSL add B edges without querying.
% Our method adds B queries (some find edges, some don't).
% Who wins?
\emph{[Results from GSL comparison]}

\subsection{Architecture Comparison}

% GCN, GAT, GraphSAGE as downstream evaluators
\emph{[Results from experiment 4e]}


\section{Analysis}
\label{sec:analysis}

\subsection{Failure Mode Evolution}
% Stacked area chart showing structural → 0 as budget spent
\emph{[Figure from fig\_results.py]}

\subsection{Submodularity}
% Marginal gains plot, submodularity ratio
\emph{[Results from theory\_submodularity.py]}

\subsection{What Kinds of Edges Matter?}
% Bridge edges vs within-cluster
\emph{[Analysis to be added]}


\section{Limitations and Future Work}
\label{sec:limitations}

\begin{itemize}
    \item The oracle is assumed perfect; real-world edge verification may be noisy.
    \item The candidate pool limits recall to $\sim$35\% of hidden edges; missed edges cannot be found.
    \item Re-scoring at every step is computationally expensive for large graphs; batch acquisition and learned policies are natural extensions.
    \item Joint label--edge budgets remain unexplored: when is an edge query more valuable than a node label?
\end{itemize}


\section{Conclusion}
\label{sec:conclusion}

We introduced active edge querying---a new form of active learning on graphs that queries structure rather than labels.
Our task-aware acquisition strategy, combining VGAE link prediction with LP-based value scoring, achieves significantly higher accuracy gains per query than link-prediction heuristics and classical active-learning baselines.
The failure diagnostic reveals that the advantage comes from eliminating structural failures---reconnecting isolated targets to labelled neighbourhoods.
We believe this opens a new direction in graph-based active learning: asking not just ``which node should I label?'' but ``which edge should I query?''


\bibliographystyle{plain}
\bibliography{references}

\end{document}
